\documentclass[10pt,a4paper]{amsart}
\usepackage{amsmath,amssymb,amsthm,graphicx}
% Retained minimal preamble for the Stacks Project chapter flat.tex (More on Flatness).
% Verbatim excerpt of the chapter preamble: Xy-pic setup (native preamble lines 14-18) and the native
% theorem declarations with the mathematical macros (native preamble lines 154-217). Copied mechanically;
% nothing added, removed or redefined. The wrapper class and the amsmath/amssymb/amsthm/graphicx base
% packages are supplied by the assembler, so the class/externaldocument part of the native preamble is omitted.
% For commutative diagrams we use Xy-pic
\usepackage[all]{xy}

% We use 2cell for 2-commutative diagrams.
\xyoption{2cell}
% Theorem environments.
%
\theoremstyle{plain}
\newtheorem{theorem}[subsection]{Theorem}
\newtheorem{proposition}[subsection]{Proposition}
\newtheorem{lemma}[subsection]{Lemma}

\theoremstyle{definition}
\newtheorem{definition}[subsection]{Definition}
\newtheorem{example}[subsection]{Example}
\newtheorem{exercise}[subsection]{Exercise}
\newtheorem{situation}[subsection]{Situation}

\theoremstyle{remark}
\newtheorem{remark}[subsection]{Remark}
\newtheorem{remarks}[subsection]{Remarks}

\numberwithin{equation}{subsection}

% Macros
%
\def\lim{\mathop{\mathrm{lim}}\nolimits}
\def\colim{\mathop{\mathrm{colim}}\nolimits}
\def\Spec{\mathop{\mathrm{Spec}}}
\def\Hom{\mathop{\mathrm{Hom}}\nolimits}
\def\Ext{\mathop{\mathrm{Ext}}\nolimits}
\def\SheafHom{\mathop{\mathcal{H}\!\mathit{om}}\nolimits}
\def\SheafExt{\mathop{\mathcal{E}\!\mathit{xt}}\nolimits}
\def\Sch{\mathit{Sch}}
\def\Mor{\mathop{\mathrm{Mor}}\nolimits}
\def\Ob{\mathop{\mathrm{Ob}}\nolimits}
\def\Sh{\mathop{\mathit{Sh}}\nolimits}
\def\NL{\mathop{N\!L}\nolimits}
\def\CH{\mathop{\mathrm{CH}}\nolimits}
\def\proetale{{pro\text{-}\acute{e}tale}}
\def\etale{{\acute{e}tale}}
\def\QCoh{\mathit{QCoh}}
\def\Ker{\mathop{\mathrm{Ker}}}
\def\Im{\mathop{\mathrm{Im}}}
\def\Coker{\mathop{\mathrm{Coker}}}
\def\Coim{\mathop{\mathrm{Coim}}}

% Boxtimes
%
\DeclareMathSymbol{\boxtimes}{\mathbin}{AMSa}{"02}

%
% Macros for moduli stacks/spaces
%
\def\QCohstack{\mathcal{QC}\!\mathit{oh}}
\def\Cohstack{\mathcal{C}\!\mathit{oh}}
\def\Spacesstack{\mathcal{S}\!\mathit{paces}}
\def\Quotfunctor{\mathrm{Quot}}
\def\Hilbfunctor{\mathrm{Hilb}}
\def\Curvesstack{\mathcal{C}\!\mathit{urves}}
\def\Polarizedstack{\mathcal{P}\!\mathit{olarized}}
\def\Complexesstack{\mathcal{C}\!\mathit{omplexes}}
% \Pic is the operator that assigns to X its picard group, usage \Pic(X)
% \Picardstack_{X/B} denotes the Picard stack of X over B
% \Picardfunctor_{X/B} denotes the Picard functor of X over B
\def\Pic{\mathop{\mathrm{Pic}}\nolimits}
\def\Picardstack{\mathcal{P}\!\mathit{ic}}
\def\Picardfunctor{\mathrm{Pic}}
\def\Deformationcategory{\mathcal{D}\!\mathit{ef}}

% --- appended 2026-10-12 (Chief): the project own definitions of the reference/slogan/history
% display environments, copied mechanically from the shared Stacks preamble lines 8-12
% (\usepackage{verbatim} plus the three \newenvironment lines, which the chapter preamble keeps
% verbatim from the project preamble). Nothing added, removed or redefined: these environments
% carry no mathematical content and the project itself typesets them as comments. Without these
% three lines a counted statement whose own text uses them fails with Environment slogan undefined.
\usepackage{verbatim}
\newenvironment{reference}{\comment}{\endcomment}
\newenvironment{slogan}{\comment}{\endcomment}
\newenvironment{history}{\comment}{\endcomment}
\usepackage{cleveref}
\crefname{originallemma}{original lemma}{original lemmas}
\Crefname{originallemma}{Original lemma}{Original lemmas}
\makeatletter
\newlabel{more-algebra-lemma-blowup-module}{{0BBJ}{source}{}{sourcepointer-1}{}}
\newlabel{more-algebra-lemma-blowup-module@cref}{{[originallemma][0][]0BBJ}{[1][0][]source}{}{}{}}
\makeatother
\crefname{originallemma}{original lemma}{original lemmas}
\Crefname{originallemma}{Original lemma}{Original lemmas}
\makeatletter
\newlabel{algebra-lemma-blowup-in-principal}{{080U}{source}{}{sourcepointer-2}{}}
\newlabel{algebra-lemma-blowup-in-principal@cref}{{[originallemma][0][]080U}{[1][0][]source}{}{}{}}
\makeatother
\crefname{originallemma}{original lemma}{original lemmas}
\Crefname{originallemma}{Original lemma}{Original lemmas}
\makeatletter
\newlabel{lemma-fibres-irreducible-flat-projective-nonnoetherian}{{05FT}{source}{}{sourcepointer-3}{}}
\newlabel{lemma-fibres-irreducible-flat-projective-nonnoetherian@cref}{{[originallemma][0][]05FT}{[1][0][]source}{}{}{}}
\makeatother
\crefname{originallemma}{original lemma}{original lemmas}
\Crefname{originallemma}{Original lemma}{Original lemmas}
\makeatletter
\newlabel{more-algebra-lemma-fitting-ideal-generate-locally}{{07ZC}{source}{}{sourcepointer-4}{}}
\newlabel{more-algebra-lemma-fitting-ideal-generate-locally@cref}{{[originallemma][0][]07ZC}{[1][0][]source}{}{}{}}
\makeatother
\crefname{originallemma}{original lemma}{original lemmas}
\Crefname{originallemma}{Original lemma}{Original lemmas}
\makeatletter
\newlabel{algebra-lemma-affine-blowup}{{07Z3}{source}{}{sourcepointer-5}{}}
\newlabel{algebra-lemma-affine-blowup@cref}{{[originallemma][0][]07Z3}{[1][0][]source}{}{}{}}
\makeatother
\crefname{originallemma}{original lemma}{original lemmas}
\Crefname{originallemma}{Original lemma}{Original lemmas}
\makeatletter
\newlabel{algebra-lemma-blowup-add-principal}{{0BBI}{source}{}{sourcepointer-6}{}}
\newlabel{algebra-lemma-blowup-add-principal@cref}{{[originallemma][0][]0BBI}{[1][0][]source}{}{}{}}
\makeatother
\crefname{originallemma}{original lemma}{original lemmas}
\Crefname{originallemma}{Original lemma}{Original lemmas}
\makeatletter
\newlabel{more-algebra-lemma-principal-fitting-ideal}{{080Z}{source}{}{sourcepointer-7}{}}
\newlabel{more-algebra-lemma-principal-fitting-ideal@cref}{{[originallemma][0][]080Z}{[1][0][]source}{}{}{}}
\makeatother
\crefname{originallemma}{original lemma}{original lemmas}
\Crefname{originallemma}{Original lemma}{Original lemmas}
\makeatletter
\newlabel{more-algebra-lemma-fitting-ideal-finite-locally-free}{{07ZD}{source}{}{sourcepointer-8}{}}
\newlabel{more-algebra-lemma-fitting-ideal-finite-locally-free@cref}{{[originallemma][0][]07ZD}{[1][0][]source}{}{}{}}
\makeatother
\crefname{originallemma}{original lemma}{original lemmas}
\Crefname{originallemma}{Original lemma}{Original lemmas}
\makeatletter
\newlabel{lemma-base-change-universally-flat}{{05DH}{source}{}{sourcepointer-9}{}}
\newlabel{lemma-base-change-universally-flat@cref}{{[originallemma][0][]05DH}{[1][0][]source}{}{}{}}
\makeatother
\title{After a suitable principal affine blowup a module that is free of rank r at the generic points of the closed fibres becomes free of rank r locally on the blowup}
\author{The Stacks Project Authors}
\date{stacks-project:flat.tex:local-snapshot-20261011; GNU FDL 1.2 or later}
\begin{document}
\maketitle
\noindent\textit{Editorial note.} All delivered mathematical text is the author's own, copied from the pinned Stacks Project chapter file flat.tex; no statement or proof step was rewritten, reordered or removed.\par\medskip
\noindent\textit{Editorial note.} The preamble is the chapter's own retained minimal preamble (sources/stacks-flat/retained-preamble.tex), a verbatim excerpt of the native Xy-pic setup and of the native theorem declarations and mathematical macros, followed by the project own definitions of the reference, slogan and history display environments (verbatim package plus three newenvironment lines) that the retained file carries at its end and that the project itself typesets as comments. No macro was invented or redefined, and the wrapper class and the amsmath/amssymb/amsthm/graphicx base packages are supplied by the assembler.\par\medskip
\noindent\textit{Editorial note.} The counted statement and proof contain no cross reference, no citation and no included file, so the delivered body is self contained.\par\medskip
\noindent\textit{Editorial note.} Canonical location of the counted proposition: Stacks tag 0810 (sources/stacks-flat/pinned-original-source.tex line 8668).\par\medskip
\noindent\textit{Allowed standard prerequisites (editorial).}\begin{itemize}
\item The noetherian-free situation of the blowup module lemma being generalised. More on Algebra, Lemma labelled more-algebra-lemma-blowup-module, cited at proof lines 8693-8695. Supplies the shape of the construction and the model for the reduction to M finitely presented.
\item Splitting S off a free R-module and the resulting coefficients. Same chapter, Lemma labelled lemma-fibres-irreducible-flat-projective-nonnoetherian, cited at proof lines 8722-8724. Makes the ideal of coefficients of the Fitting ideal a finitely generated ideal of R.
\item Fitting ideals, their base change and their relation to freeness. More on Algebra, Lemmas labelled more-algebra-lemma-fitting-ideal-generate-locally, more-algebra-lemma-principal-fitting-ideal and more-algebra-lemma-fitting-ideal-finite-locally-free, cited at proof lines 8740-8742, 8756-8758 and 8766-8768. Converts the condition on the Fitting ideal into local generation by r elements and then into local freeness of rank r.
\item Affine blowup in one principal element and the nonzerodivisor property. Algebra, Lemmas labelled algebra-lemma-affine-blowup, algebra-lemma-blowup-add-principal and algebra-lemma-blowup-in-principal, cited at proof lines 8748-8750 and 8772-8774. Identifies I'R' with a'R', shows a' is a nonzerodivisor, and identifies R'\_a with R\_f.
\end{itemize}\medskip
\ifdefined\hypertarget\hypertarget{sourcepointer-1}{}\fi
\noindent\textit{Source reference.} Original lemma 0BBJ: sources/stacks-chow/tags.txt, source lines 14247--14247. Cross-chapter standard prerequisite located in the project tag list staged in the local snapshot; invoked at proof line 8695 of the counted proof. It is retained only as a navigation pointer; every load-bearing step is inside the delivered span.\par\medskip
\ifdefined\hypertarget\hypertarget{sourcepointer-2}{}\fi
\noindent\textit{Source reference.} Original lemma 080U: sources/stacks-chow/tags.txt, source lines 10194--10194. Cross-chapter standard prerequisite located in the project tag list staged in the local snapshot; invoked at proof line 8709 of the counted proof. It is retained only as a navigation pointer; every load-bearing step is inside the delivered span.\par\medskip
\ifdefined\hypertarget\hypertarget{sourcepointer-3}{}\fi
\noindent\textit{Source reference.} Original lemma 05FT: sources/stacks-flat/pinned-original-source.tex, source lines 2247--2247. Same-chapter pointer into the author environment lemma at source line 2247; invoked at proof line 8717 of the counted proof. It is retained only as a navigation pointer; every load-bearing step is inside the delivered span.\par\medskip
\ifdefined\hypertarget\hypertarget{sourcepointer-4}{}\fi
\noindent\textit{Source reference.} Original lemma 07ZC: sources/stacks-chow/tags.txt, source lines 10142--10142. Cross-chapter standard prerequisite located in the project tag list staged in the local snapshot; invoked at proof line 8733 of the counted proof. It is retained only as a navigation pointer; every load-bearing step is inside the delivered span.\par\medskip
\ifdefined\hypertarget\hypertarget{sourcepointer-5}{}\fi
\noindent\textit{Source reference.} Original lemma 07Z3: sources/stacks-chow/tags.txt, source lines 10133--10133. Cross-chapter standard prerequisite located in the project tag list staged in the local snapshot; invoked at proof line 8741 of the counted proof. It is retained only as a navigation pointer; every load-bearing step is inside the delivered span.\par\medskip
\ifdefined\hypertarget\hypertarget{sourcepointer-6}{}\fi
\noindent\textit{Source reference.} Original lemma 0BBI: sources/stacks-chow/tags.txt, source lines 14246--14246. Cross-chapter standard prerequisite located in the project tag list staged in the local snapshot; invoked at proof line 8742 of the counted proof. It is retained only as a navigation pointer; every load-bearing step is inside the delivered span.\par\medskip
\ifdefined\hypertarget\hypertarget{sourcepointer-7}{}\fi
\noindent\textit{Source reference.} Original lemma 080Z: sources/stacks-chow/tags.txt, source lines 10199--10199. Cross-chapter standard prerequisite located in the project tag list staged in the local snapshot; invoked at proof line 8756 of the counted proof. It is retained only as a navigation pointer; every load-bearing step is inside the delivered span.\par\medskip
\ifdefined\hypertarget\hypertarget{sourcepointer-8}{}\fi
\noindent\textit{Source reference.} Original lemma 07ZD: sources/stacks-chow/tags.txt, source lines 10143--10143. Cross-chapter standard prerequisite located in the project tag list staged in the local snapshot; invoked at proof line 8768 of the counted proof. It is retained only as a navigation pointer; every load-bearing step is inside the delivered span.\par\medskip
\ifdefined\hypertarget\hypertarget{sourcepointer-9}{}\fi
\noindent\textit{Source reference.} Original lemma 05DH: sources/stacks-flat/pinned-original-source.tex, source lines 1907--1907. Same-chapter pointer into the author environment lemma at source line 1907; invoked at proof line 8776 of the counted proof. It is retained only as a navigation pointer; every load-bearing step is inside the delivered span.\par\medskip
\begin{lemma}
\label{lemma-helper-blowup-affine-space}
Let $R$ be a ring and let $f \in R$. Let $r\geq 0$ be an integer.
Let $R \to S$ be a ring map and let $M$ be an $S$-module. Assume
\begin{enumerate}
\item $R \to S$ is of finite presentation and flat,
\item every fibre ring $S \otimes_R \kappa(\mathfrak p)$ is
geometrically integral over $R$,
\item $M$ is a finite $S$-module,
\item $M_f$ is a finitely presented $S_f$-module,
\item for all $\mathfrak p \in R$, $f \not \in \mathfrak p$ with
$\mathfrak q = \mathfrak pS$ the module $M_{\mathfrak q}$ is free
of rank $r$ over $S_\mathfrak q$.
\end{enumerate}
Then there exists a finitely generated ideal $I \subset R$ with
$V(f) = V(I)$ such that for all $a \in I$ with $R' = R[\frac{I}{a}]$
the quotient
$$
M' = (M \otimes_R R')/a\text{-power torsion}
$$
over $S' = S \otimes_R R'$ satisfies the following: for every prime
$\mathfrak p' \subset R'$ there exists a $g \in S'$,
$g \not \in \mathfrak p'S'$ such that $M'_g$ is a free $S'_g$-module
of rank $r$.
\end{lemma}

\begin{proof}
This lemma is a generalization of
More on Algebra, Lemma \ref{more-algebra-lemma-blowup-module};
we urge the reader to read that proof first.
Choose a surjection $S^{\oplus n} \to M$, which is possible by (1).
Choose a finite submodule $K \subset \Ker(S^{\oplus n} \to M)$
such that $S^{\oplus n}/K \to M$ becomes an isomorphism after inverting $f$.
This is possible by (4). Set $M_1 = S^{\oplus n}/K$ and suppose we can
prove the lemma for $M_1$. Say $I \subset R$ is the corresponding ideal.
Then for $a \in I$ the map
$$
M_1' = (M_1 \otimes_R R')/a\text{-power torsion}
\longrightarrow
M' = (M \otimes_R R')/a\text{-power torsion}
$$
is surjective. It is also an isomorphism after inverting $a$ in $R'$
as $R'_a = R_f$, see Algebra, Lemma \ref{algebra-lemma-blowup-in-principal}.
But $a$ is a nonzerodivisor on $M'_1$, whence the displayed map is an
isomorphism. Thus it suffices to prove the lemma in case $M$ is a finitely
presented $S$-module.

\medskip\noindent
Assume $M$ is a finitely presented $S$-module satisfying (3).
Then $J = \text{Fit}_r(M) \subset S$ is a finitely generated ideal.
By Lemma \ref{lemma-fibres-irreducible-flat-projective-nonnoetherian}
we can write $S$ as a direct summand of a free
$R$-module: $\bigoplus_{\alpha \in A} R = S \oplus C$.
For any element $h \in S$ writing $h = \sum a_\alpha$ in the
decomposition above, we say that the $a_\alpha$ are the coefficients of $h$.
Let $I' \subset R$ be the ideal of coefficients
of elements of $J$. Multiplication by an element of $S$ defines
an $R$-linear map $S \to S$, hence $I'$ is generated by the coefficients
of the generators of $J$, i.e., $I'$ is a finitely generated ideal.
We claim that $I = fI'$ works.

\medskip\noindent
We first check that $V(f) = V(I)$. The inclusion $V(f) \subset V(I)$ is
clear. Conversely, if $f \not \in \mathfrak p$, then
$\mathfrak q =  \mathfrak p S$ is not an element of $V(J)$ by
property (5) and More on Algebra, Lemma
\ref{more-algebra-lemma-fitting-ideal-generate-locally}.
Hence there is an
element of $J$ which does not map to zero in $S \otimes_R \kappa(\mathfrak p)$.
Thus there exists an element of $I'$ which is not contained in
$\mathfrak p$, so $\mathfrak p \not \in V(fI') = V(I)$.

\medskip\noindent
Let $a \in I$ and set $R' = R[\frac{I}{a}]$. We may write $a = fa'$
for some $a' \in I'$. By Algebra, Lemmas \ref{algebra-lemma-affine-blowup} and
\ref{algebra-lemma-blowup-add-principal} we see that $I' R' = a'R'$
and $a'$ is a nonzerodivisor in $R'$. Set $S' = S \otimes_S R'$.
Every element $g$ of $JS' = \text{Fit}_r(M \otimes_S S')$ can be
written as $g = \sum_\alpha c_\alpha$ for some $c_\alpha \in I'R'$.
Since $I'R' = a'R'$ we can write $c_\alpha = a'c'_\alpha$ for some
$c'_\alpha \in R'$ and $g = (\sum c'_\alpha)a' = g' a'$ in $S'$.
Moreover, there is an $g_0 \in J$ such that $a' = c_\alpha$
for some $\alpha$. For this element we have $g_0 = g'_0 a'$ in $S'$
where $g'_0$ is a unit in $S'$.
Let $\mathfrak p' \subset R'$ be
a prime ideal and $\mathfrak q' = \mathfrak p'S'$.
By the above we see that $JS'_{\mathfrak q'}$ is the
principal ideal generated by the nonzerodivisor $a'$.
It follows from More on Algebra, Lemma
\ref{more-algebra-lemma-principal-fitting-ideal}
that $M'_{\mathfrak q'}$ can be generated by $r$ elements.
Since $M'$ is finite, there exist $m_1, \ldots, m_r \in M'$ and
$g \in S'$, $g \not \in \mathfrak q'$ such that the corresponding map
$(S')^{\oplus r} \to M'$ becomes surjective after inverting $g$.

\medskip\noindent
Finally, consider the ideal $J' = \text{Fit}_{r - 1}(M')$.
Note that $J'S'_g$ is generated by the coefficients of relations between
$m_1, \ldots, m_r$ (compatibility of Fitting ideal with base change).
Thus it suffices to show that $J' = 0$, see
More on Algebra, Lemma
\ref{more-algebra-lemma-fitting-ideal-finite-locally-free}.
Since $R'_a = R_f$ (Algebra, Lemma \ref{algebra-lemma-blowup-in-principal})
and $M'_a = M_f$ we see from (5)
that $J'_a$ maps to zero in $S_{\mathfrak q''}$ for any prime
$\mathfrak q'' \subset S'$ of the form $\mathfrak q'' = \mathfrak p''S'$
where $\mathfrak p'' \subset R'_a$. Since
$S'_a \subset \prod_{\mathfrak q''\text{ as above}} S'_{\mathfrak q''}$
(as $(S'_a)_{\mathfrak p''} \subset S'_{\mathfrak q''}$ by
Lemma \ref{lemma-base-change-universally-flat})
we see that $J'R'_a = 0$. Since $a$ is a nonzerodivisor in $R'$ we
conclude that $J' = 0$ and we win.
\end{proof}

\end{document}
